% =====================================================================
%  PROJECT RECORD — analytical, reportage and visualisation work on the
%  Indian downstream aluminium tariff study.
%  A handover record and a source for CV / portfolio framing.
% =====================================================================
\documentclass[11pt, a4paper]{article}
\usepackage[a4paper, top=2.0cm, bottom=2.0cm, left=2.1cm, right=2.1cm]{geometry}
\usepackage{fontspec}
\IfFontExistsTF{DejaVu Sans}
  {\setmainfont{DejaVu Sans}}
  {\IfFontExistsTF{Verdana}{\setmainfont{Verdana}}{\setmainfont{Arial}}}
\usepackage{booktabs}
\usepackage{amsmath}
\usepackage{enumitem}
\usepackage{xcolor}
\usepackage{titlesec}
\usepackage{float}
\usepackage{tcolorbox}
\usepackage{longtable}
\emergencystretch=2.5em
\hbadness=10000

\definecolor{shade}{gray}{0.945}
\newtcolorbox{glance}{colback=shade, colframe=black!55, boxrule=0.5pt,
  arc=2pt, left=7pt, right=7pt, top=5pt, bottom=5pt}

\titleformat{\section}{\large\bfseries}{\thesection}{0.7em}{}
\titleformat{\subsection}{\normalsize\bfseries}{\thesubsection}{0.7em}{}
\setlength{\parskip}{0.35em}

\title{\textbf{Project Record}\\
\large Rationalizing Tariff Structures in the Indian Downstream Aluminium Industry\\
\normalsize Analytical work, reportage and visualisation}
\author{}
\date{}

\begin{document}
\maketitle
\vspace{-0.9cm}

\begin{glance}
\textbf{At a glance.} An evidence-based policy submission arguing for rationalisation
of the basic customs duty on primary aluminium (HS~7601). The work comprised: an
effective-rate-of-protection analysis across thirteen product lines and six trading
partners; an industrial-statistics analysis built from 3,705 observations of the
Annual Survey of Industries; a public-finance analysis of the tariff-versus-GST
trade-off; and the production of a 45-page report with 21 figures, supported by an
automated verification harness and three technical companion documents.
\end{glance}

\section{Objective and scope}

The study set out to answer four questions: whether India's present tariff structure
places downstream aluminium products at a measurable disadvantage; what that
disadvantage costs in value addition and employment; whether the fiscal objection to
correcting it holds; and what policy change follows. Scope was fixed at
HS~7601--7616, the chapter-76 lines spanning unwrought metal through fabricated
articles.

The central analytical problem was that the tariff evidence and the industrial
evidence had to describe \emph{the same industry}. Resolving that --- deriving the
industrial classification from the tariff schedule rather than from convention ---
was the methodological core of the project.

\section{Data sources}

\begin{table}[H]
\centering\small
\begin{tabular}{p{4.6cm}p{4.0cm}p{6.4cm}}
\toprule
\textbf{Source} & \textbf{Coverage} & \textbf{Used for} \\
\midrule
Annual Survey of Industries, All-India summary & 3,705 observations; 15 indicators, 19 NIC classes, 13 years (2011--12 to 2023--24) & All industrial statistics: value added, employment, wages, workers, contract labour, net profit \\
\addlinespace
MOSPI Supply--Use Tables & 26 supply and use files, 2011--12 to 2023--24 & Effective protection estimates; input-cost shares \\
\addlinespace
Tariff schedules & 256 tariff lines, HS 7601--7616, six partners & MFN and preferential rates; the ERP basket \\
\addlinespace
Effective corporate tax rates & By year, 2011--12 to 2023--24 & Corporate-tax base estimation \\
\addlinespace
Ministry of Mines, OECD, GTRI, CEA, Union Budget & Published & Capacity, demand, trade composition, public capital expenditure \\
\bottomrule
\end{tabular}
\end{table}

\section{Analytical work}

\subsection{Segmentation: deriving the classification from the tariff schedule}

The prior draft partitioned the ASI on convention. The work replaced this with a
rule derived from Corden's definition of an \emph{activity} --- a stage of conversion
from traded input to traded output --- so that a NIC class enters the downstream
aggregate if and only if its principal output falls within the thirteen ERP lines
HS~7604--7616. This produced upstream $=$ NIC 2420 and downstream $=$ NIC 2432,
2511, 2599, and required excluding four classes previously included whose outputs
fall in HS Chapters 85 and 88.

The rule was cross-checked against the OECD's aluminium value-chain segmentation and
against the material-flow literature, and the divergence from the latter was tested
rather than argued: reassigning the semi-fabrication class moves the headline result
from 3.14 to 2.77 and changes no conclusion.

\subsection{Series construction}

Thirteen annual series were built for each segment and each constituent class:
gross value added, output, inputs, wages, persons engaged, total workers, contract
workers, net profit, and estimated corporate tax. From these, nine derived measures:
value added per person engaged, labour-intensity ratio, wage share, contract share,
input-cost share, employment yield per unit of tax, index series, three-year endpoint
growth, and year-on-year log volatility.

\subsection{Principal findings}

\begin{table}[H]
\centering\small
\begin{tabular}{p{6.6cm}p{4.0cm}p{4.4cm}}
\toprule
\textbf{Finding} & \textbf{Result} & \textbf{Robustness} \\
\midrule
Labour intensity: employment sustained per unit of value added & \textbf{3.14$\times$} downstream; range 3.08--6.18, median 3.93 & Holds in 12 of 12 years; exceeds the Ministry of Mines' 2.5$\times$ benchmark in all \\
\addlinespace
Labour share of value added & \textbf{38.4\%} vs \textbf{17.6\%} & Series never cross in 12 years; uses only monetary variables \\
\addlinespace
Employment per unit of estimated corporate tax & \textbf{151} vs \textbf{39} persons per crore & Ordering holds in 12 of 12 years \\
\addlinespace
Structural products (NIC 2511, HS 7610) & $+$\textbf{15.0\%} nominal over 11 years; workforce flat & The single most disprotected ERP line \\
\addlinespace
Volatility of value added & Upstream \textbf{3.6$\times$} downstream & sd of year-on-year log growth \\
\addlinespace
Input-cost amplification & \textbf{4.91$\times$} & Derived from a measured input share of 83.1\% \\
\bottomrule
\end{tabular}
\end{table}

\subsection{Analytical results established rather than asserted}

Three results were derived formally rather than argued or simulated, and these are
the most durable output of the project.

\begin{enumerate}[leftmargin=1.6em, itemsep=4pt]
\item \textbf{Apportionment invariance.} The objection that the NIC classes are not
exclusively aluminium was closed by proof rather than sensitivity runs. An
apportioned segment productivity is necessarily a weighted average of its class
productivities, hence bounded by their range for \emph{any} positive apportionment
vector; the upstream segment being a single class, its coefficient cancels. The
labour-intensity ratio therefore lies in $[3.01,\ 3.38]$ whatever apportionment is
assumed. Verified against 200,000 random draws.

\item \textbf{The per-tonne decomposition.} The employment-per-tonne ratio factors
exactly into the per-rupee ratio times value added per tonne
($3.14 \times 1.81 = 5.7$). This reconciled the report's own measure with the much
larger ratios circulating in industry discussion, and showed those to arise from
comparing mismatched employment universes.

\item \textbf{The transfer-to-revenue identity.} In the tariff-versus-GST analysis,
the ratio of surplus transferred to revenue collected reduces to
$a/c = \phi(1-\mu)/[\mu(1-f)]$ --- import penetration, pass-through and the duty-free
share. Price, exchange rate and both elasticities cancel. The result is therefore
robust to the parameters least well known.
\end{enumerate}

\subsection{Public-finance analysis}

A separate analysis addressed the fiscal objection to rationalisation: partial
equilibrium welfare accounting (producer surplus, revenue, and both deadweight
triangles, using supply and demand elasticities from the metals literature), the
revenue-replacement arithmetic against the GST base, and an induced-expansion offset
derived from the amplification factor. Headline results: \textbf{INR 17.5}
transferred per rupee collected, \textbf{54 paise} of national income destroyed per
rupee raised, and a revenue requirement replaceable by \textbf{0.18 of a percentage
point} of GST.

\section{Academic grounding}

The argument was grounded in primary sources rather than secondary summary, with
27 references in the final bibliography. The load-bearing citations:

\begin{table}[H]
\centering\small
\begin{tabular}{p{5.4cm}p{9.6cm}}
\toprule
\textbf{Source} & \textbf{Role in the argument} \\
\midrule
Corden (1966) & Defines the activity as the unit of effective-protection analysis; supplies the segmentation rule \\
Balassa (1965, 1968) & Precedent for negative effective protection; tariff escalation \\
Pathania \& Bhattacharjea (2020) & Modified Corden estimator; Indian inverted duty structures \\
Diamond \& Mirrlees (1971) & Production efficiency: intermediate goods should not be taxed \\
Bhagwati \& Ramaswami (1963) & Target the distortion at source; trade instruments are second-best \\
Emran \& Stiglitz (2005); Keen (2008) & The informality objection to VAT-for-tariff substitution, and its limits \\
Baunsgaard \& Keen (2010) & Empirical evidence on revenue recovery after liberalisation \\
Bohlin \& Widell (2006); Miller \& Blair (2009) & Technology assumptions; factor-content invariance \\
Cullen \& Allwood (2013); Bertram et al.\ (2009) & Material-flow taxonomy of the aluminium system \\
Chaurey (2015); Goldberg et al.\ (2010) & Contract labour; imported intermediate inputs and product growth \\
OECD (2019); Viner (1950) & Value-chain segmentation; trade diversion \\
\bottomrule
\end{tabular}
\end{table}

\section{Visualisation}

Twenty-one figures across the report, seven of them rebuilt for the industrial
statistics section. All are vector charts generated directly from the computed
series, with the plotting coordinates emitted by the analysis code so that a figure
cannot drift from its data.

\subsection{Design constraints and how they were met}

\begin{enumerate}[leftmargin=1.6em, itemsep=3pt]
\item \textbf{Monochrome legibility.} The report is a policy submission that will be
printed and photocopied. Every series is distinguished by three redundant cues ---
line style, marker shape and marker fill --- so that no series depends on colour or
grey level. Bar charts pair a hatched light fill against a solid dark one.

\item \textbf{No text merging with the plot.} Every printed value sits on an opaque
white chip, so a plot line or gridline passing beneath cannot obscure it. Where two
series run close or cross, per-point labelling was replaced by direct annotation of
endpoints and turning points.

\item \textbf{Legends outside the plotting area}, never overlying data.

\item \textbf{Values on the chart.} No number in a figure has to be read off an axis.

\item \textbf{Captions carry a source line only.} All construction, derivation and
domain detail was moved to the appendix, so that the narrative reads without
interruption.
\end{enumerate}

\subsection{Editorial decisions on the figure set}

The section was reduced from eight figures to seven single-panel figures, cutting
Chapter 4 from roughly seventeen pages to nine. Three charts were removed on the
principle that a figure earns its place only if it advances the argument: a
value-added-per-person chart whose ratio was the following figure; a persons-engaged
index panel superseded by the labour-intensity chart; and a corporate-tax
\emph{levels} panel, on which the comparison runs against the report's own case ---
retained as a sentence of text so the concession is made, but not given a chart.

\section{Quality assurance}

\subsection{The verification harness}

An automated checker recomputes every series from source, parses every plotting
coordinate out of the document, and tests whether each plotted number is a correct
rounding of the computed value \emph{at the precision printed}. It also re-checks the
statistics quoted in prose. Current status: \textbf{all checks pass across 158
plotted points and 30 quoted statistics}.

\subsection{Errors found and corrected}

The harness caught four errors that had survived several drafts: three plotted points
rounded twice (136.47 printed as 137, 214.47 as 215, 139.47 as 140), and an
amplification factor computed from a rounded input share and printed as 4.92 instead
of 4.91. A separate manual check found a caption claiming NIC 2511 spent eight of
eleven years below its base level, where the correct count is seven. None changed a
conclusion; none would have been found by reading.

\subsection{Analytical corrections made during review}

Two substantive errors were caught and corrected in the analysis itself: a robustness
claim that conflated leave-one-out with reassignment of a class (3.13 stated where
2.77 was correct), and a transfer-to-revenue ratio computed on downstream fabrication
throughput where the correct base is all domestic primary metal sold domestically
(4.7:1 stated where 17.5:1 was correct).

\subsection{Standards of claim}

Causal language was audited and withdrawn where it had not been earned. The report
now states explicitly that its evidence is association rather than an identified
causal effect, names the absent counterfactual and the uncontrolled construction
cycle, and removes a timing claim that twelve annual observations cannot support.
Known weak points --- an assumed parameter, an industry-estimated tonnage, a tax
proxy, a frame that excludes unincorporated units --- are stated in the documents
rather than left to be discovered.

\section{Deliverables}

\begin{table}[H]
\centering\small
\begin{tabular}{p{5.6cm}p{1.6cm}p{7.8cm}}
\toprule
\textbf{Document} & \textbf{Extent} & \textbf{Purpose} \\
\midrule
Main report & 45 pp, 21 figures & The policy submission \\
Author's notes & 22 pp & Every calculation behind every figure, with worked arithmetic; variable and coefficient glossary; the arguments and the figure carrying each \\
Tariff-versus-GST note & 8 pp & Public-finance analysis; declared baseline, welfare accounting, sensitivity \\
Decision record & 2 pp & Reviewer comments: what was adopted, what was not, and why \\
\midrule
\multicolumn{3}{l}{\textit{Analysis and build code}} \\
Series computation & --- & Builds all series; emits plotting coordinates \\
Verification harness & --- & Independent re-derivation and rounding audit \\
Report assembly & --- & Deterministic build with cross-reference and bibliography injection \\
Public-finance arithmetic & --- & Declared baseline, results, sensitivity \\
\bottomrule
\end{tabular}
\end{table}

\section{Summary lines, for a CV or capability statement}

\begin{itemize}[leftmargin=1.4em, itemsep=4pt]
\item Built the industrial-statistics evidence base for a policy submission on
aluminium tariff reform, constructing thirteen-year series for nine measures from
3,705 ASI observations and reconciling the industrial classification to the tariff
schedule so that the statistical and trade evidence described the same industry.

\item Established the central finding --- that value added in downstream fabrication
sustains 3.1 times the employment of value added in primary metal --- and demonstrated
it held in every year of the record, corroborating an official estimate from an
independent source.

\item Closed the principal methodological objection analytically rather than by
simulation, proving the headline ratio bounded within $[3.01, 3.38]$ for any
apportionment of the underlying classes, and verifying the bound against 200,000
draws.

\item Conducted a public-finance analysis of the tariff-versus-GST trade-off,
deriving a transfer-to-revenue identity independent of price, exchange rate and
elasticity, and quantifying the revenue at issue as replaceable by 0.18 of a
percentage point of GST.

\item Designed and produced 21 monochrome-legible vector figures under print
constraints, with plotting coordinates generated directly from the analysis so that
figures could not drift from their data.

\item Built an automated verification harness that re-derives every plotted point and
quoted statistic from source; it caught four numerical errors that had survived
several drafts.

\item Grounded the argument in primary academic sources across trade theory, optimal
taxation and industrial ecology, and audited the report's causal language, withdrawing
claims the data could not support.
\end{itemize}

\end{document}
